% © 2026 Ing. David Jaroš — CC BY-NC-ND 4.0
% Licensed under CC BY-NC-ND 4.0. See LICENSE.md.
\section{The biquaternionic frame connection $\Omega_\mu$}
\label{sec:canonical:biquaternion_connection}

A connection specifies how local Lorentz/biquaternionic frames are compared at
neighbouring spacetime points.  The UBT tetrad is
\begin{equation}
 E_\mu=\mathcal N_0^{-1/2}D_\mu\Theta.
\end{equation}
The action of $D_\mu$ must always be stated explicitly; left, right, adjoint,
and two-sided actions are not interchangeable.

\subsection{Classical Lorentz connection reconstructed from tetrad and torsion}
Write
$E_\mu=i e_\mu{}^0\mathbf1+e_\mu{}^k\mathbf e_k$.  Let
\begin{equation}
 T^a=de^a+\omega^a{}_b\wedge e^b
 =\frac12T^a{}_{bc}e^b\wedge e^c.
\end{equation}
For every nondegenerate tetrad and specified torsion, there is exactly one
metric-compatible Lorentz connection,
\begin{equation}
 \boxed{\omega^a{}_b=\mathring\omega^a{}_b(e)+K^a{}_b(T),}
 \label{eq:canonical:omega-e-T}
\end{equation}
where
\begin{equation}
 \mathring\omega_\mu{}^a{}_b
 =e_b{}^\nu\left(
 \mathring\Gamma^\rho{}_{\mu\nu}(g)e_\rho{}^a
 -\partial_\mu e_\nu{}^a
 \right)
 \label{eq:canonical:omega-from-e}
\end{equation}
and, with the displayed torsion convention,
\begin{equation}
 \boxed{K_{abc}=\frac12(T_{cab}-T_{abc}-T_{bca}).}
 \label{eq:canonical:contorsion}
\end{equation}
Thus the connection has no residual kinematic ambiguity once $(e,T)$ are
fixed.  In the torsion-free classical GR branch, $T=0$, $K=0$, and
$\omega=\mathring\omega(e)$.

The biquaternionic/spin connection is the lift
\begin{equation}
 \Omega_\mu=\rho_{\mathrm{spin}*}(\omega_\mu)
 =\frac14\omega_{\mu ab}\Sigma^{ab}
\end{equation}
for the chosen representation.  The complete proof is in
\texttt{canonical/gr\_closure/gap\_10omega\_connection\_elimination.tex}.

\subsection{Relation to Christoffel symbols}
The coordinate and frame descriptions are joined by
\begin{equation}
 \partial_\mu e_\nu{}^a
 -\Gamma^\rho{}_{\mu\nu}e_\rho{}^a
 +\omega_\mu{}^a{}_b e_\nu{}^b=0.
 \label{eq:canonical:tetrad-postulate}
\end{equation}
$\Gamma^\rho{}_{\mu\nu}$ transports coordinate indices, while $\omega$ and
$\Omega$ transport local Lorentz/biquaternionic indices.  They encode the same
geometry in different representations and are not related by
$\Gamma=\operatorname{Re}\Omega$.

\subsection{Gauge covariance and Lorentz-slice preservation}
Under $e\mapsto\Lambda e$,
\begin{equation}
 \omega\mapsto\Lambda\omega\Lambda^{-1}-d\Lambda\,\Lambda^{-1}.
\end{equation}
This is a local Lorentz gauge change, not a new physical degree of freedom.
Since every metric-compatible $\omega_\mu$ lies in $\mathfrak{so}(1,3)$, its
parallel transport preserves $\eta_{ab}$ and the real Lorentz slice, including
branches with contorsion.  Hence GAP-10L-CONN is closed.  GAP-10L-DYN remains
open because the complete $\Theta$ equations and sources must also preserve the
slice.

\subsection{Representation selection on $\Theta$}
A naive one-sided action
\begin{equation}
 D^L_\mu\Theta=\partial_\mu\Theta+A_\mu\Theta
\end{equation}
obeys $[D^L_\mu,D^L_\nu]\Theta=F^L_{\mu\nu}\Theta$.  Under torsion-free
tetrad compatibility, invertible $\Theta$ would then force $F^L=0$.  Therefore
this one-sided regular representation is not the generic curved-GR route.

The algebra-native viable form is the two-sided bimodule derivative
\begin{equation}
 \boxed{
 D_\mu\Theta=\partial_\mu\Theta+A_\mu\Theta-\Theta B_\mu,}
 \label{eq:canonical:two-sided-D}
\end{equation}
for which
\begin{equation}
 \boxed{
 [D_\mu,D_\nu]\Theta
 =F^A_{\mu\nu}\Theta-\Theta F^B_{\mu\nu}.}
 \label{eq:canonical:two-sided-curvature}
\end{equation}
Torsion-free antisymmetric compatibility requires
$F^A_{\mu\nu}\Theta=\Theta F^B_{\mu\nu}$; if $\Theta$ is invertible, the two
curvatures are conjugate rather than zero.  This removes the one-sided
flatness obstruction but does not yet derive the precise relation between
$A_\mu$ and $B_\mu$ from the action.

\subsection{Flat limit and explicit representer}
In a flat inertial torsion-free sector,
\begin{equation}
 \Gamma^\rho{}_{\mu\nu}=0,
 \qquad \Omega_\mu=0,
 \qquad D_\mu=\partial_\mu.
\end{equation}
The explicit field
\begin{equation}
 \Theta_{\mathrm{SR}}=\Theta_0+\sqrt{\mathcal N_0}
 (i x^0\mathbf1+x^k\mathbf e_k)
\end{equation}
generates the Minkowski tetrad and has vanishing second spacetime derivatives.

\subsection{Full UBT status}
The kinematic connection-reconstruction problem is closed:
\begin{itemize}
 \item GAP-10$\Omega$-KIN: specified $(e,T)$ uniquely determine $\omega$;
 \item GAP-10$\Omega$-GR: $T=0$ gives the Levi--Civita spin connection;
 \item GAP-10L-CONN: metric-compatible transport preserves the Lorentz slice.
\end{itemize}
Open dynamical questions are:
\begin{itemize}
 \item GAP-10T-DYN: derive $T$ from the UBT action and sources;
 \item fix the exact paired left/right connection and involution on $\Theta$;
 \item GAP-10I-CURVED: solve the implicit curved system
       $E_\mu=D_\mu\Theta/\sqrt{\mathcal N_0}$;
 \item GAP-10L-DYN: prove slice preservation by the full dynamics.
\end{itemize}
